\clearpage\pagestyle{otMain}\setlength{\otBlocEspacement}{2.5mm}\fontsize{11}{15}\selectfont\corrigesChapitre{19}{Données et enquêtes}
\begin{blocLivre}{cor-19-1}\begin{corrige}{1}{Le bon outil}\demarche{Se demander d’où vient l’information.}
1. Un $\res{\text{questionnaire}}$. 2. Un $\res{\text{comptage}}$, sur place. 3. Des $\res{\text{mesures}}$, avec une toise et une unité commune.
\end{corrige}\end{blocLivre}
\begin{blocLivre}{cor-19-2}\begin{corrige}{2}{Construire un tableau}\demarche{Compter chaque réponse, puis contrôler le total.}
Chat : $\res8$ ; chien : $\res4$ ; aucun : $\res4$ ; poisson : $\res2$ ; lapin : $\res2$. Total : $8+4+4+2+2=20$, le compte est bon.
\end{corrige}\end{blocLivre}
\begin{blocLivre}{cor-19-3}\begin{corrige}{3}{Lire un diagramme en barres}\demarche{Lire chaque hauteur sur l’axe gradué.}
Le $\res{\text{mercredi}}$, avec 35 livres. Total : $15+20+35+25+30=\res{125}$. Écart : $35-15=\res{20}$ livres.
\label{fig-cor-19-3}\figureGeo[corrige]{don-3}
\end{corrige}\end{blocLivre}
\begin{blocLivre}{cor-19-4}\begin{corrige}{4}{Lire un diagramme circulaire}\demarche{Associer chaque secteur à une fraction du disque.}
Vidéos : la moitié, $\res{2\text{ h}}$. Jeux : un quart, $\res{1\text{ h}}$. Messages : un quart, $\res{1\text{ h}}$.
\label{fig-cor-19-4}\figureGeo[corrige]{don-4}
\end{corrige}\end{blocLivre}
\begin{blocLivre}{cor-19-5}\begin{corrige}{5}{Lire une courbe}\demarche{Lire l’abscisse et l’ordonnée du point le plus haut.}
Maximum : $\res{21\ ^\circ\mathrm{C}}$ à $\res{15\text{ h}}$. Hausse : $21-8=\res{13\ ^\circ\mathrm{C}}$.
\label{fig-cor-19-5}\figureGeo[corrige]{don-5}
\end{corrige}\end{blocLivre}
\begin{blocLivre}{cor-19-6}\begin{corrige}{6}{Relevé de mesures}\demarche{Calculer les écarts d’une semaine à la suivante.}
\begin{center}\renewcommand{\arraystretch}{1.35}\begin{tabular}{c|ccccc}\toprule
Semaine & 1 & 2 & 3 & 4 & 5\\\midrule
Taille (cm) & 2 & 5 & 9 & 14,5 & 17\\
\bottomrule\end{tabular}\end{center}Croissances : 3 cm ; 4 cm ; $5{,}5$ cm ; $2{,}5$ cm. Il a le plus grandi entre les semaines $\res{3\text{ et }4}$.
\end{corrige}\end{blocLivre}
\begin{blocLivre}{cor-19-7}\begin{corrige}{7}{Tableau à double entrée}\demarche{Additionner en ligne et en colonne, puis contrôler.}
Totaux des lignes : 26 et 25. Totaux des colonnes : 18 ; 10 ; 23. À pied : $\res{18}$ élèves. En tout : $\res{51}$ élèves, et $26+25=18+10+23=51$.
\end{corrige}\end{blocLivre}
\begin{blocLivre}{cor-19-8}\begin{corrige}{8}{Filtrer pour choisir}\demarche{Appliquer les critères l’un après l’autre.}
Mercredi : échecs, théâtre, chorale, escalade. Prix au plus 10 € : échecs, chorale, escalade. Places restantes : $\res{\text{échecs et escalade}}$.
\end{corrige}\end{blocLivre}
\begin{blocLivre}{cor-19-9}\begin{corrige}{9}{Un graphique trompeur}\demarche{Vérifier où commence l’axe vertical.}
L’axe vertical commence à 100 : la seconde barre paraît 4 fois plus haute, alors que les ventes n’ont augmenté que de $108-102=\res6$ articles, soit environ $\res{6\text{~\%}}$. Avec un axe partant de 0, les barres seraient presque égales.
\label{fig-cor-19-9}\figureGeo[corrige]{don-9}
\end{corrige}\end{blocLivre}
\begin{blocLivre}{cor-19-10}\begin{corrige}{10}{Construire un diagramme}\demarche{Choisir l’échelle, puis tracer chaque barre.}
Les barres, en vert, montent à 12, 9, 5 et 4. La barre de la danse s’arrête entre les graduations 8 et 10.
\label{fig-cor-19-10}\figureGeo[corrige]{don-10}
\end{corrige}\end{blocLivre}
\begin{blocLivre}{cor-19-11}\begin{corrige}{11}{Le gaspillage à la cantine}\demarche{Additionner, puis appliquer un pourcentage.}
Total : $3{,}2+2{,}8+3{,}5+2{,}5=\res{12\ \mathrm{kg}}$. Une baisse de 25~\%, c’est 3 kg de moins : objectif $\res{9\ \mathrm{kg}}$ par semaine.
\end{corrige}\end{blocLivre}
\begin{blocLivre}{cor-19-12}\begin{corrige}{12}{Les insectes du jardin}\demarche{Filtrer selon le nombre de jours non nuls.}
Jours avec observation : abeille 4, coccinelle 2, papillon 4, fourmi 5, sauterelle 1. Au moins 3 jours : $\res{\text{abeille, papillon et fourmi}}$. La plus comptée : la fourmi, avec $\res{54}$ observations.
\end{corrige}\end{blocLivre}
\begin{blocLivre}{cor-19-13}\begin{corrige}{13}{Le climat d’une ville}\demarche{Lire chaque représentation avec son propre axe.}
Le plus chaud : $\res{\text{juillet}}$ (23 °C). Le plus sec : $\res{\text{juillet}}$ (18 mm). Le plus pluvieux : $\res{\text{novembre}}$ (70 mm). En été, il fait chaud et il pleut peu.
\label{fig-cor-19-13}\figureGeo[corrige]{don-13}
\end{corrige}\end{blocLivre}
\begin{blocLivre}{cor-19-14}\begin{corrige}{14}{Planifier une enquête}\demarche{Suivre les six étapes du cours.}
Population : les élèves de la classe. Recueil : un questionnaire où chacun note sa durée, en minutes. Tableau : trois colonnes, $\res{\text{moins de 10 min, de 10 à 20 min, plus de 20 min}}$, avec les effectifs et un total. Puis un diagramme en barres et une phrase de conclusion.
\end{corrige}\end{blocLivre}
\begin{blocLivre}{cor-19-15}\begin{corrige}{15}{Choisir un séjour}\demarche{Filtrer critère par critère.}
Au bord de la mer : Bretagne, Vendée, Corse, Landes. À 400 € au plus : Vendée et Landes. Au moins 5 jours : seule la $\res{\text{Vendée}}$ convient.
\end{corrige}\end{blocLivre}
